% Ad Familiares 10.19 --- headnote
% ad-familiares-10-19, around 26 May 43 BC, Rome

\textit{Cicero to L. Munatius Plancus, written at Rome around
26 May 43 BC --- Perseus dateline \textit{Scr. Romae circ.
vii K. Iun. a. 711 (43)}. A short reply to a now-lost
dispatch of Plancus which had reached the Senate and made a
strong impression there. Cicero opens with the private note
--- he had not been looking for thanks, but the expression
of affection he received was unexpectedly pleasant --- and
then turns at once to the business: the Senate received
Plancus's letter, for its substance and for the weight of
its language, ``wonderfully well.''}

\textit{The second section is the pressing peroration in
miniature that the spring of 43 keeps generating from
Cicero's pen: put your shoulder to it, finish the war,
crush Antony. The note that ``my affection for the
fatherland is hardly greater than for your glory'' is the
characteristic Ciceronian way of binding ambition to duty
in his correspondents: the man who finishes Antony off
will, by the same act, have served the commonwealth and
made his name immortal. Within days, however, the news
from Gaul will undo everything: Lepidus's army will go
over to Antony, and Plancus will retreat back across the
Is\`ere.}
